% appendix.tex
\clearpage
\appendix

\section{Hyperparameters}
\label{app:hyperparams}

Table~\ref{tab:hyperparams} lists all hyperparameters used across
experimental configurations.

\begin{table}[!htb]
\centering
\caption{Complete hyperparameter specification.}
\label{tab:hyperparams}
\footnotesize
\begin{tabular}{llll}
\toprule
Component & Parameter & Value & Notes \\
\midrule
\multicolumn{4}{l}{\emph{Embedding}} \\
& Model & all-MiniLM-L6-v2 & 82M params, 384-dim \\
& Normalization & L2 & cosine similarity \\
\midrule
\multicolumn{4}{l}{\emph{Chunking (fixed)}} \\
& Chunk sizes & \{256, 512, 1024\} & tokens \\
& Overlap & 10\% & 26, 51, 102 tokens \\
& Separators & \texttt{["\textbackslash n\textbackslash n", "\textbackslash n", ". ", " ", ""]} & recursive \\
\midrule
\multicolumn{4}{l}{\emph{Retrieval}} \\
& Top-$k$ & \{3, 5\} & cosine similarity \\
& Vector store & ChromaDB v0.4 & persistent \\
\midrule
\multicolumn{4}{l}{\emph{Generation}} \\
& Models & Mistral-7B, Llama-3-8B & via Ollama v0.1.29 \\
& Temperature & 0.1 & low stochasticity \\
\midrule
\multicolumn{4}{l}{\emph{Cross-encoder}} \\
& Model & ms-marco-MiniLM-L-6-v2 & 66M params \\
& Output & raw logits & not probabilities \\
\midrule
\multicolumn{4}{l}{\emph{NLI scorer}} \\
& Model & nli-deberta-v3-base & 435M params \\
& Threshold & 0.50 & hallucination cutoff \\
\midrule
\multicolumn{4}{l}{\emph{HCPC-v2}} \\
& $\tau_\text{sim}$ & 0.45 & cosine threshold \\
& $\tau_\text{ce}$ & $-$0.20 & CE logit threshold \\
& $N_\text{prot}$ & 2 & rank protection \\
& $N_\text{max}$ & 2 & max refine per query \\
& Sub-chunk size & 256 tokens & for splitting \\
& Sub-overlap & 32 tokens & \\
& Merge max & 4096 chars & $\approx$1024 tokens \\
\midrule
\multicolumn{4}{l}{\emph{Reranker}} \\
& Fetch-$k$ & 7 & over-fetch depth \\
& Return-$k$ & 3 & final context size \\
\bottomrule
\end{tabular}
\end{table}

\FloatBarrier
\section{Chunk size ablation and variance analysis}
\label{app:ablation}

Tables~\ref{tab:ablation_squad} and~\ref{tab:ablation_pubmedqa} report
the full 12-configuration ablation over chunk size, top-$k$, and prompt
strategy. All experiments use Mistral-7B with $n = 30$ queries per
configuration.

\begin{table}[!htb]
\centering
\caption{Full chunk size ablation on SQuAD (Mistral-7B, $n = 30$).}
\label{tab:ablation_squad}
\small
\begin{tabular}{cclcc}
\toprule
Chunk & $k$ & Prompt & Faithfulness & Halluc. \\
\midrule
256  & 3 & strict & 0.6732 & 10.0\% \\
256  & 3 & cot    & 0.6835 & 6.7\% \\
256  & 5 & strict & 0.7254 & 10.0\% \\
256  & 5 & cot    & 0.7253 & 10.0\% \\
\midrule
512  & 3 & strict & 0.7827 & 0.0\% \\
512  & 3 & cot    & 0.7781 & 0.0\% \\
512  & 5 & strict & 0.7861 & 3.3\% \\
512  & 5 & cot    & 0.7737 & 6.7\% \\
\midrule
1024 & 3 & strict & 0.8001 & 3.3\% \\
1024 & 3 & cot    & 0.7891 & 0.0\% \\
1024 & 5 & strict & 0.8274 & 3.3\% \\
1024 & 5 & cot    & 0.8254 & 3.3\% \\
\bottomrule
\end{tabular}
\end{table}

\begin{table}[!htb]
\centering
\caption{Full chunk size ablation on PubMedQA (Mistral-7B, $n = 30$).}
\label{tab:ablation_pubmedqa}
\small
\begin{tabular}{cclcc}
\toprule
Chunk & $k$ & Prompt & Faithfulness & Halluc. \\
\midrule
256  & 3 & strict & 0.5591 & 30.0\% \\
256  & 3 & cot    & 0.5802 & 20.0\% \\
256  & 5 & strict & 0.5585 & 30.0\% \\
256  & 5 & cot    & 0.5636 & 26.7\% \\
\midrule
512  & 3 & strict & 0.6028 & 16.7\% \\
512  & 3 & cot    & 0.6196 & 13.3\% \\
512  & 5 & strict & 0.5844 & 20.0\% \\
512  & 5 & cot    & 0.5906 & 13.3\% \\
\midrule
1024 & 3 & strict & 0.5963 & 20.0\% \\
1024 & 3 & cot    & 0.6167 & 10.0\% \\
1024 & 5 & strict & 0.5849 & 16.7\% \\
1024 & 5 & cot    & 0.5881 & 13.3\% \\
\bottomrule
\end{tabular}
\end{table}

\paragraph{Chunk-size summary.} Mean faithfulness by chunk size on
SQuAD: $256 = 0.7018$, $512 = 0.7801$, $1024 = 0.8105$. Mean
hallucination: $256 = 9.17\%$, $512 = 2.5\%$, $1024 = 2.5\%$. On
PubMedQA: faithfulness $256 = 0.5654$, $512 = 0.5994$, $1024 = 0.5965$;
hallucination $256 = 26.67\%$, $512 = 15.83\%$, $1024 = 15.0\%$. The
$256 \to 1024$ improvement is large on SQuAD (10.87~pp faithfulness,
6.67~pp hallucination reduction) and moderate on PubMedQA (3.11~pp
faithfulness, 11.67~pp hallucination reduction).

\FloatBarrier
\subsection{ANOVA results}
\label{app:anova}

Table~\ref{tab:anova} reports one-way ANOVA results for each factor.

\begin{table}[!htb]
\centering
\caption{One-way ANOVA for faithfulness by factor.}
\label{tab:anova}
\small
\begin{tabular}{llcccc}
\toprule
Dataset & Factor & $F$ & $p$ & $\eta^2$ & Sig. \\
\midrule
SQuAD    & Chunk size      & 19.720 & $<$0.001 & 0.100 & Yes \\
SQuAD    & Prompt strategy & 0.046  & 0.830    & $<$0.001 & No \\
SQuAD    & Top-$k$         & 2.918  & 0.089    & 0.012 & No \\
\midrule
PubMedQA & Chunk size      & 3.664  & 0.027    & 0.020 & Yes \\
PubMedQA & Prompt strategy & 1.122  & 0.290    & 0.005 & No \\
PubMedQA & Top-$k$         & 2.323  & 0.128    & 0.010 & No \\
\bottomrule
\end{tabular}
\end{table}

Chunk size is the only factor reaching significance on both datasets.
On SQuAD, $F = 19.72$ with $p < 0.001$; on PubMedQA, $F = 3.66$ with
$p = 0.027$. Prompt strategy is non-significant on both datasets
($p = 0.830$ on SQuAD, $p = 0.290$ on PubMedQA), and the effect size
is negligible ($\eta^2 < 0.005$).

\FloatBarrier
\subsection{Cohen's \texorpdfstring{$d$}{d} effect sizes}
\label{app:cohend}

Table~\ref{tab:cohend} reports pairwise Cohen's $d$ between factor levels.

\begin{table}[!htb]
\centering
\caption{Cohen's $d$ effect sizes for faithfulness.}
\label{tab:cohend}
\small
\begin{tabular}{lcc}
\toprule
Comparison & SQuAD & PubMedQA \\
\midrule
\multicolumn{3}{l}{\emph{Chunk size}} \\
256 vs 512  & 0.82 & 0.31 \\
256 vs 1024 & 1.41 & 0.29 \\
512 vs 1024 & 0.55 & 0.03 \\
\midrule
\multicolumn{3}{l}{\emph{Prompt strategy}} \\
strict vs CoT & 0.02 & 0.11 \\
\midrule
\multicolumn{3}{l}{\emph{Top-$k$}} \\
$k = 3$ vs $k = 5$ & 0.18 & 0.16 \\
\bottomrule
\end{tabular}
\end{table}

The 256-to-1024 chunk size effect on SQuAD ($d = 1.41$) is large by
conventional standards ($d > 0.8$). The 512-to-1024 effect ($d = 0.55$)
is medium. On PubMedQA, all chunk size effects are small ($d < 0.35$),
reflecting the domain mismatch that limits the benefit of larger chunks.
Prompt strategy effects are negligible on both datasets ($d \leq 0.11$).

\FloatBarrier
\subsection{Variance decomposition by model}
\label{app:variance_model}

Table~\ref{tab:variance_model} partitions the variance in faithfulness
among the three configuration factors, computed separately for each
model-dataset combination. This uses the multi-model data from the
extended experiments (Llama-3-8B and Mistral-7B, $n = 100$ per condition).

\begin{table}[!htb]
\centering
\caption{Variance decomposition: proportion of faithfulness variance
explained by each factor, by model and dataset.}
\label{tab:variance_model}
\small
\begin{tabular}{llccc}
\toprule
Model & Dataset & Chunk size & Top-$k$ & Prompt \\
\midrule
\multicolumn{2}{l}{\emph{Overall}} \\
--- & SQuAD    & 90.1\% & 9.8\%  & 0.2\% \\
--- & PubMedQA & 61.2\% & 26.1\% & 12.7\% \\
\midrule
\multicolumn{2}{l}{\emph{By model}} \\
Mistral-7B  & SQuAD    & 72.4\% & 9.0\%  & 18.6\% \\
Llama-3-8B  & SQuAD    & 53.3\% & 25.2\% & 21.5\% \\
Mistral-7B  & PubMedQA & 1.5\%  & 97.9\% & 0.6\% \\
Llama-3-8B  & PubMedQA & 95.2\% & 0.4\%  & 4.4\% \\
\bottomrule
\end{tabular}
\end{table}

The model-level breakdown reveals substantial heterogeneity. On SQuAD,
chunk size dominates for both models (72.4\% for Mistral, 53.3\% for
Llama-3), consistent with the overall 90.1\% figure. On PubMedQA, the
models diverge sharply: Mistral's faithfulness is governed by top-$k$
(97.9\%), while Llama-3's is governed by chunk size (95.2\%). This
suggests that optimal RAG configurations are model-specific and should
not be assumed to transfer between architectures.

\paragraph{Why chunk size matters.} Chunk size affects faithfulness
through three reinforcing mechanisms. First, larger chunks contain
complete sentences and paragraphs; a 256-token chunk may truncate
mid-thought, producing lower NLI entailment scores even for qualitatively
correct answers. Second, with $k = 3$ and chunk = 1024, the generator
receives approximately 3,000 tokens of context versus 750 at chunk = 256;
more context provides redundant evidence that reduces hallucination
risk. Third, larger chunks preserve the logical flow of source
documents---topic sentences, supporting details, and transitions appear
together, maintaining the coherence that smaller chunks fragment.

The 256-to-1024 effect on SQuAD ($d = 1.41$) reflects all three
mechanisms acting together. The 512-to-1024 effect ($d = 0.55$) is
smaller because 512-token chunks already capture most complete sentences
and moderate redundancy.

\paragraph{Why prompt strategy does not matter.} The maximum
faithfulness difference between strict and chain-of-thought prompts on
SQuAD, holding chunk size and top-$k$ constant, is $0.0121$
(at chunk~$=1024$, $k=3$: strict~$=0.8001$, CoT~$=0.7891$). Both prompts
instruct the model to answer from the context. When the context is
coherent and contains the answer, both produce faithful responses; when
the context is fragmented or the answer is absent, neither can prevent
hallucination. For RAG faithfulness, prompt engineering yields returns
roughly $450\times$ smaller than chunk-size selection
($\eta^2_{\text{chunk}} / \eta^2_{\text{prompt}} = 0.901 / 0.002$).

\FloatBarrier
\section{Re-ranking by chunk size}
\label{app:reranking}

Table~\ref{tab:reranker_chunk} reports re-ranking results across chunk
sizes (Llama-3-8B, $n = 100$ on SQuAD, $n = 50$ on PubMedQA).

\begin{table}[!htb]
\centering
\caption{Re-ranking effect by chunk size. $\Delta$ columns show
reranked minus baseline.}
\label{tab:reranker_chunk}
\small
\begin{tabular}{llcccc}
\toprule
Dataset & Chunk & Base Faith. & Rerank Faith. & $\Delta$ Faith. & $\Delta$ Halluc. \\
\midrule
SQuAD    & 256  & 0.7233 & 0.7265 & +0.003 & $-$3.0 pp \\
SQuAD    & 512  & 0.7523 & 0.7544 & +0.002 & $-$3.0 pp \\
SQuAD    & 1024 & 0.7739 & 0.7619 & $-$0.012 & 0.0 pp \\
\midrule
PubMedQA & 256  & 0.6136 & 0.6455 & +0.032 & $-$10.0 pp \\
PubMedQA & 512  & 0.5862 & 0.5957 & +0.010 & +6.0 pp \\
PubMedQA & 1024 & 0.5853 & 0.5806 & $-$0.005 & $-$2.0 pp \\
\bottomrule
\end{tabular}
\end{table}

Re-ranking interacts with chunk size in a non-obvious way. At small chunk
sizes (256), re-ranking consistently helps: it improves faithfulness and
reduces hallucination on both datasets, with the largest effect on
PubMedQA ($-$10 pp hallucination). At large chunk sizes (1024),
re-ranking slightly \emph{hurts} faithfulness on both datasets. At
intermediate sizes (512 on PubMedQA), re-ranking improves faithfulness
but \emph{increases} hallucination by 6 pp.

The pattern suggests that re-ranking is most useful when the initial
retrieval is noisy (small chunks produce more candidates of variable
quality) and least useful when chunks are already large enough to be
self-contained. At chunk = 1024, the cross-encoder's preference for
keyword-matching relevance can override the embedding model's broader
semantic ranking, replacing a contextually appropriate passage with one
that matches query terms more precisely but contributes less to overall
coherence.

\FloatBarrier
\section{Multi-model chunk size comparison}
\label{app:multimodel}

Table~\ref{tab:multimodel} reports faithfulness and hallucination rates
for both models across chunk sizes (extended experiments, $n = 100$ per
condition on SQuAD).

\begin{table}[!htb]
\centering
\caption{Multi-model comparison by chunk size.}
\label{tab:multimodel}
\small
\begin{tabular}{llccc}
\toprule
Model & Dataset & Chunk 256 & Chunk 512 & Chunk 1024 \\
\midrule
\multicolumn{5}{l}{\emph{Faithfulness}} \\
Llama-3  & SQuAD    & 0.7184 & 0.7530 & 0.7591 \\
Mistral  & SQuAD    & 0.7248 & 0.7548 & 0.7676 \\
Llama-3  & PubMedQA & 0.6151 & 0.5851 & 0.5868 \\
Mistral  & PubMedQA & 0.5798 & 0.5797 & 0.5818 \\
\midrule
\multicolumn{5}{l}{\emph{Hallucination rate}} \\
Llama-3  & SQuAD    & 4.5\%  & 1.75\% & 1.0\% \\
Mistral  & SQuAD    & 8.25\% & 4.5\%  & 3.75\% \\
Llama-3  & PubMedQA & 16.0\% & 19.0\% & 18.5\% \\
Mistral  & PubMedQA & 26.5\% & 28.5\% & 25.5\% \\
\bottomrule
\end{tabular}
\end{table}

Both models show the same chunk-size trend on SQuAD: larger chunks
produce higher faithfulness and lower hallucination. On PubMedQA, the
trend is weaker and non-monotonic for Llama-3 (faithfulness peaks at
chunk = 256), consistent with the domain mismatch limiting the benefit
of more context when that context is poorly matched.

\FloatBarrier
\section{Threshold sweep details}
\label{app:sweep}

The threshold sweep tested 168 fixed configurations
(7 $\tau_\text{sim}$ values $\times$ 6 $\tau_\text{ce}$ values
$\times$ 4 $N_\text{prot}$ values) plus 4 adaptive configurations
per dataset. Evaluation used the proxy score:

\begin{equation}
  \text{proxy} = 0.4 \cdot \text{CCS} - 0.3 \cdot \text{var}(S)
               - 0.2 \cdot H(S) + 0.1 \cdot \bar{q}
\end{equation}

\noindent where $\text{var}(S)$ is the embedding variance,
$H(S)$ is the retrieval entropy, and $\bar{q}$ is the mean
query-chunk similarity.

On SQuAD, the top-10 fixed configurations all achieve CCS = 1.0
with proxy score = 0.2587, compared to the default configuration's
0.2425. The optimal region is broad: configurations with
$\tau_\text{sim} \geq 0.35$ and $N_\text{prot} = 0$ produce
identical proxy scores because the AND gate rarely fires when
the embedding model matches the corpus.

On PubMedQA, the top-10 configurations achieve proxy score = 0.2820
(default: 0.2496, $\Delta$ = +13\%). These concentrate at
$\tau_\text{sim} \in [0.50, 0.60]$ with $N_\text{prot} = 0$, where
more chunks are eligible for refinement and the proxy rewards higher
coherence after refinement.

Adaptive (percentile-based) thresholds underperform fixed thresholds on
both datasets. The best adaptive configuration matches the default fixed
configuration (proxy $= 0.2425$ on SQuAD, $0.2496$ on PubMedQA). Fixed
thresholds, once tuned to the dataset, are sufficient; adaptive
adjustment adds complexity without benefit.

\FloatBarrier
\section{Per-query failure analysis}
\label{app:failure}

We manually examined hallucinated responses from HCPC-v1 on SQuAD
($n = 10$). Three failure patterns emerged:

\begin{enumerate}
  \item \textbf{Evidence truncation} (6/10): The refined sub-chunk did
        not contain the sentence needed to answer the question. The
        model generated a plausible but unsupported answer using its
        parametric memory.

  \item \textbf{Context gap bridging} (3/10): The model inferred a
        connection between two non-adjacent sub-chunks that was not
        supported by either. The gap left by refinement was filled
        with fabricated reasoning.

  \item \textbf{Over-completion} (1/10): The model extended its answer
        beyond what the context supported, appending details from
        parametric memory after correctly answering the initial question.
\end{enumerate}

Under HCPC-v2, the only failure mode observed (3/3 hallucinated responses
in extended experiments) was evidence truncation, and all three cases
occurred on queries where the rank protection gate was active. This
suggests that selective refinement occasionally prevents refinement of
a passage that genuinely needs it, when the weak chunk happens to rank in
the top 2.

\FloatBarrier
\section{Latency measurements}
\label{app:latency}

\begin{table}[!htb]
\centering
\caption{Mean generation latency (seconds, $\pm$ std dev).}
\label{tab:latency}
\small
\begin{tabular}{llccc}
\toprule
Model & Dataset & Chunk = 256 & Chunk = 512 & Chunk = 1024 \\
\midrule
Llama-3  & SQuAD    & 2.20 $\pm$ 1.16 & 3.23 $\pm$ 1.01 & 4.19 $\pm$ 1.59 \\
Llama-3  & PubMedQA & 6.77 $\pm$ 43.3 & 5.87 $\pm$ 1.70 & 8.71 $\pm$ 2.78 \\
Mistral  & SQuAD    & 3.41 $\pm$ 1.78 & 4.56 $\pm$ 1.68 & 9.04 $\pm$ 91.3 \\
Mistral  & PubMedQA & 6.37 $\pm$ 2.15 & 10.3 $\pm$ 3.28 & 13.0 $\pm$ 3.80 \\
\bottomrule
\end{tabular}
\end{table}

Latency scales roughly linearly with chunk size (more context tokens
to process). The high standard deviations on some cells (e.g., Llama-3
PubMedQA chunk = 256: $\pm$43.3s) reflect occasional outlier queries
where the model produced unusually long responses.

The HCPC-v2 refinement step adds 2--3 seconds per query (cross-encoder
scoring + sub-splitting). On SQuAD, where refinement triggers for 16.7\%
of queries, the amortized overhead is approximately 0.4 seconds per query.
On PubMedQA, where 73.3\% of queries trigger refinement, the overhead is
approximately 1.8 seconds per query.

\FloatBarrier
\section{Implementation details}
\label{app:implementation}

The pipeline is implemented in Python 3.11. Key dependencies:

\begin{itemize}
  \item \texttt{langchain} and \texttt{langchain-community} for
        text splitting and Ollama integration
  \item \texttt{sentence-transformers} for embedding and cross-encoder
        inference
  \item \texttt{chromadb} v0.4 for vector storage with cosine distance
  \item \texttt{torch} with MPS backend for Apple Silicon acceleration
  \item \texttt{scipy} for ANOVA and correlation tests
  \item \texttt{pandas} for result aggregation
\end{itemize}

The HCPC-v2 retriever (\texttt{src/hcpc\_v2\_retriever.py}, 550 lines)
is the most complex component. It maintains per-query diagnostic logs
that record: whether refinement was triggered, the number of chunks
before and after refinement, the best cosine and cross-encoder scores
before and after, the CCS, approximate token count, threshold mode,
and effective threshold values.

All experiments use the same 30 (or 100, depending on phase) queries per
condition, sampled randomly from the validation split at the start of
each experiment and held fixed across all conditions within that
experiment. This ensures that condition comparisons are not confounded
by query difficulty variation.

\section{Real-world case studies}
\label{app:case_studies}

To make the refinement paradox concrete beyond the aggregate metrics,
we present a small case-study set: queries on which the same generator
gives a faithful answer with raw retrieval, hallucinates after
per-passage refinement, and recovers under the coherence-gated policy.
The cases are auto-selected from the per-query logs by the script
\texttt{experiments/build\_qualitative\_example.py} as the queries
with the largest baseline$\to$HCPC-v$1$ faithfulness gap.

% Auto-generated by experiments/build_qualitative_example.py (--top mode)
\begin{figure*}[!htb]\centering
\medskip\noindent\fbox{\begin{minipage}{0.96\linewidth}\small \textbf{Case 1.} \textit{Dataset:} squad \quad \textit{Model:} mistral \\
\textbf{Question.}~What city did Super Bowl 50 take place in? \\
\textbf{Ground truth.}~\textit{Santa Clara}

\textbf{Baseline} ($f=0.942$): The Super Bowl 50 took place in Santa Clara, which is part of the San Francisco Bay Area. \\ 
\textbf{HCPC-v1} ($f=0.478$): San Francisco Bay Area \\ 
\textbf{HCPC-v2} ($f=0.942$, recovered): The Super Bowl 50 took place in Santa Clara, which is part of the San Francisco Bay Area.\end{minipage}}


\medskip\noindent\fbox{\begin{minipage}{0.96\linewidth}\small \textbf{Case 2.} \textit{Dataset:} squad \quad \textit{Model:} mistral \\
\textbf{Question.}~What city did Super Bowl 50 take place in? \\
\textbf{Ground truth.}~\textit{Santa Clara}

\textbf{Baseline} ($f=0.942$): The Super Bowl 50 took place in Santa Clara, which is part of the San Francisco Bay Area. \\ 
\textbf{HCPC-v1} ($f=0.478$): San Francisco Bay Area \\ 
\textbf{HCPC-v2} ($f=0.942$, recovered): The Super Bowl 50 took place in Santa Clara, which is part of the San Francisco Bay Area.\end{minipage}}


\medskip\noindent\fbox{\begin{minipage}{0.96\linewidth}\small \textbf{Case 3.} \textit{Dataset:} squad \quad \textit{Model:} mistral \\
\textbf{Question.}~What year did the Denver Broncos secure a Super Bowl title for the third time? \\
\textbf{Ground truth.}~\textit{2015}

\textbf{Baseline} ($f=0.795$): The provided context does not specify the exact year when the Denver Broncos secured a Super Bowl title for the third time. However, it is implied that this event happened in Super Bowl 50, which was held in 2016. \\ 
\textbf{HCPC-v1} ($f=0.485$): I cannot find this information in the provided context. The context only mentions that the Denver Broncos won the AFC Championship Game and secured home-field advantage throughout the AFC playoffs, but it does not speci\textbackslash\{\}dots \\ 
\textbf{HCPC-v2} ($f=0.795$, recovered): The provided context does not specify the exact year when the Denver Broncos secured a Super Bowl title for the third time. However, it is implied that this event happened in Super Bowl 50, which was held in 2016.\end{minipage}}


\medskip\noindent\fbox{\begin{minipage}{0.96\linewidth}\small \textbf{Case 4.} \textit{Dataset:} squad \quad \textit{Model:} llama-3.3-70b \\
\textbf{Question.}~If Roman numerals were used, what would Super Bowl 50 have been called? \\
\textbf{Ground truth.}~\textit{Super Bowl L}

\textbf{Baseline} ($f=0.502$): If Roman numerals were used, Super Bowl 50 would have been called Super Bowl L. \\ 
\textbf{HCPC-v1} ($f=0.491$): Super Bowl L. \\ 
\textbf{HCPC-v2} ($f=0.502$, recovered): If Roman numerals were used, Super Bowl 50 would have been called Super Bowl L.\end{minipage}}

\caption{\textbf{Real-world case studies of the refinement paradox.} Five queries where the same question is answered correctly with raw retrieval, incorrectly after per-passage refinement, and correctly again under coherence-gated refinement. Faithfulness $f$ is the NLI entailment score (higher is better; $\geq 0.5$ is the conventional faithful threshold). All cases satisfy $\textrm{baseline} = \textrm{faithful}$, $\textrm{HCPC-v1} = \textrm{hallucinated}$, $\textrm{HCPC-v2} = \textrm{faithful}$.}
\label{fig:qualitative_cases}
\end{figure*}

